\documentclass[../../main.tex]{subfiles}  % Two levels up to main.tex
\begin{document}

\setupappendix{A}  % See `usercommands.tex` for definitions

\section{Chapter 3 supplementary material}\label{sec:app-a}

\subsection{Supplemental information}\label{sec:s2-info}

Appendices in this template are typeset single-spaced by default.
This can be changed in \texttt{backmatter.tex}, though consult with formatting review and the guidelines for your department. 
There's a variety of different ways you can structure an appendix based on your field. 
In science, these might be the supplementary materials that accompany individual chapters or published manuscripts. 
In other fields, they may correspond to derivations, proofs, algorithms, equations, methods, analyses, protocols, or attempted methodologies not germane to the main body text. 

\subsection{Supplemental figures}\label{sec:s2-figs}

Notice that all figures in this template are emumerated first by the chapter number, followed by a second number representing the order in which it appears in the chapter. 
For Appendices, this template correctly enumerates its encompassed figures and tables by its letter rather than an additional chapter number (see \Cref{fig:a1-f1}). 
The same logic applies for subpanel references, such as \subpanelrange{fig:a1a}{fig:a1b}{B}.

\begin{figure}[ht]
    \centering
    \includegraphics[width=14cm]{./Plots/a1_sample.pdf}
    {
    \phantomsubcaption\label{fig:a1a}
    \phantomsubcaption\label{fig:a1b}
    }
    \caption{
        Appendix A example figure. 
        \subref{fig:a1a} Shows raster plots of single neuron activity relative to the onset of target (orange) and non-target (teal) stimuli, followed by a \acrfull{psth} of their trial averages.
        \subref{fig:a1b} Polar histogram of preferred local field potential phases for recorded neurons following target and non-target stimuli.  
    }
    \label{fig:a1-f1}
\end{figure}

\subsection{Supplemental tables}

\end{document}